\section{数据预处理}

\subsection{数据来源与总体特征}

六份附件数据体量大、维度多、时间跨度长，必须预处理数据。预处理分成检查、构造、探索三个环节。先做完整性检查与异常检测，疑似异常的观测逐条甄别并标记。确认数据干净后，再构造 GPU-hour、时延可达矩阵等派生变量，容量折算与迁移决策共用一套计量规则。最后做探索性分析，把对建模有参考价值的数据事实挑出，供各子问题建模时直接参考。

本题数据由六份附件构成。\texttt{workload\_trace.xlsx} 记录 50\,000 条 AI 任务的完整信息，每条含类型、到达与最晚完成小时、GPU 需求量、执行时长、最大时延和来源区域等关键属性。\texttt{region\_time\_data.xlsx} 给出六区域 2\,407 个逐时断面，共 $6\times2407=14\,442$ 行，覆盖逐时节点电价 234--1\,096\,元/MWh、碳强度 0.196--0.688\,tCO$_2$/MWh、可用新能源 500--1\,100\,MW 与 GPU 利用率，功率平衡残差最大仅 0.0002\,MW。\texttt{GPU\_information.xlsx} 给出 GPU 总量 600--1\,600、保留比例 8\%--10\% 与 PUE 1.25--1.38。\texttt{storage\_information.xlsx} 记录储能容量 300--900\,MWh、初始 SOC 45\%、最小 SOC 10\%，最大购电 340--550\,MW、最大售电 0--220\,MW，A/B/C 区无外送能力，售电上限恒为 0。\texttt{network\_latency.xlsx} 给出 36 对单向时延，同区 5\,ms，东部区域间 12--15\,ms，东西部 58--82\,ms，西部区域间 5--25\,ms。\texttt{power\_mapping.xlsx} 规定训练、批量、实时三类等效 IT 功率分别为 0.16、0.10、0.08\,MW/GPU。

\subsection{数据清洗与异常处理}

六份附件均无缺失值。50\,000条任务记录的字段完备，TaskID没有重复，记录中到达小时覆盖0--2399。14\,442行给出$2407\times6$的完整信息。四份参数表均完整填值且数据一致。

异常检测共甄别出三类特别的数据特征，处理方式各有依据。\texttt{region\_time\_data.xlsx}的GPU\_Utilization\_Percent列有44条超过100\%的观测值（最大135.58\%，集中在E/F区，占0.3\%），经核实源于赛题基准运行方案下少数时段的负载瞬时尖峰，故予以保留并标记作为基准参考，建模时按GPU-hour重叠折算实际利用率并用于容量约束，以避免任务并发超出GPU可用量。弃电量均值约537\,MW，占可用新能源约67\%，将用作新能源利用率提升空间的衡量基准。A/B/C区GridSell恒为0的50\%行属于无外送能力的真实设定，因此不做填补。

\subsection{关键派生变量构造}

为支撑后续各子问题的建模与分析，我们从原始数据中构造了如下四类关键派生变量。

\begin{enumerate}[label=\arabic*.]
    \item \textbf{GPU-hour。}GPU-hour是容量约束折算与负载统计的基础计量单位，定义为单任务GPU需求量与执行时长的乘积，即 $\text{GH}_i = \text{GPU\_Demand}_i \times \text{EstimatedDuration\_min}_i / 60$。容量约束据此要求任意小时同时运行任务的GPU需求量之和不超过总可用GPU数量，任务每占用一小时即计入相应GPU需求。全部50\,000个任务按精确小数计算的累计GPU-hour，即代表全部算力需求的总量。

    \item \textbf{网络时延矩阵。}网络时延矩阵用于确定任务迁移的可达性。我们将$6\times6$单向时延表与三类任务的时延上限（MaxLatency）交叉匹配，得到各自的可达区域对。实时推理仅限$\le20$\,ms的路径，故只能在A/B/C之间（12--15\,ms）或E/F之间（18\,ms）互迁，而D区没有迁移目的地，其至E/F的时延均超过20\,ms。批量推理排除超过80\,ms的路径，A$\to$F（82\,ms）因此被禁止，A$\to$E（78\,ms）仍可使用；AI训练则36对区域组合全部可达。

    \item \textbf{PUE$\times$均价排序键。}该排序键用于刻画跨区真实用电成本差异，我们将PUE与区域均价相乘，得到每MWh IT负载真实用电成本的排序键$K_r=\text{PUE}_r\times\overline{\text{Price}}_r$，供迁移决策快速比较跨区边际成本。排序键由低到高为E$<$F$<$D$<$C$<$B$<$A，与西部低成本高能效、东部高成本的格局相符。

    \item \textbf{任务候选数。}任务候选数汇总上述可达性结果，量化每个任务的可迁移范围。我们对每个任务按来源与时延约束搜索可达目的地：训练任务全图可达，均有6个候选；批量推理平均有4.85个候选（34/36区域对可达）；实时推理的东部来源任务有3个候选（A/B/C互迁），西部来源任务则只能在E/F之间互迁。没有可行迁移目的地的任务约占1.0\%。
\end{enumerate}

\subsection{探索性数据分析与关键事实}

我们围绕任务结构、区域差异与电力价格展开工作后，分析出可供建模参考的关键事实。

从GPU-hour结构看，AI训练任务是算力负荷的主要来源。训练任务数量虽只占约三分之一（16\,559条），但单任务GPU需求高、执行时长长，功率高，时长中位数约240\,min，等效IT功率0.16\,MW/GPU，是实时推理的2倍，三者叠加使其贡献约80\%的GPU-hour总量，是碳感知调度最主要的调节对象。与之相对，实时推理对时延最敏感、难以跨区调度，但需求总量少。

六区域在GPU容量、电价与碳强度上均存在明显差异。GPU容量以D区1\,600为最大，C区只有600。PUE存在东西差异，西部E/F为1.25最低，东部A区1.38最高，能效差约10\%。电价梯度主要落在东西部之间，E区均价约374、F区约393、A区约707\,元/MWh，价差达1.2--1.9倍。碳强度东西之比为2.8倍，A区为0.668、E区为0.237\,tCO$_2$/MWh。

逐时可用新能源处于500--1\,100\,MW，各区域各小时都有非零值，但基准方案下弃电率高达约67\%。弃电的主要原因是A/B/C区没有外送通道，富余新能源无法跨区消纳。六区域储能总容量约3\,300\,MWh。储能初始SOC统一为45\%，最小SOC为10\%，充放电效率约92\%--94\%。

\subsection{预处理小结}

最后，50\,000条任务轨迹全部保留，14\,442行逐时断面经异常标记后完整可用，并且GPU-hour、时延可达矩阵、PUE$\times$均价排序键与任务候选数等派生变量也已构造。接下来我们将在对应章节说明各子问题的特殊处理。
